

% Palette (print-safe, clear hierarchy for dementia consultation system)
\definecolor{fillCore}{RGB}{210,235,220}     % light green: core contribution
\definecolor{fillPartial}{RGB}{232,245,236}  % lighter green: partial contribution
\definecolor{fillContext}{RGB}{235,235,235}  % grey: context / not contributed
\definecolor{fillExisting}{RGB}{245,245,245} % very light grey: pre-existing / setting
\definecolor{accentGreen}{RGB}{70,140,95}    % green text/accent

\begin{figure}[p]
\centering
\begin{tikzpicture}[
  font=\small,
  node distance=7mm,
  box/.style={draw=black, rounded corners, align=center, inner sep=6pt, text width=0.90\textwidth},
  arrow/.style={-{Latex[length=2mm]}, thick, draw=black}
]

% Pre-existing / setting context (very light grey)
\node[box, fill=fillExisting] (patient) {
\textbf{Clinical consultation (primary care)} \\
Patient presenting with memory concerns or cognitive decline \\
Complex interactions involving patient and caregiver emotions
};

% Core contributions (light green fill)
\node[box, fill=fillCore, below=of patient] (navigation) {
\textbf{Navigation Map guidance (Objective 3)} \\
Three-phase structured workflow: Recognition, Evaluation, Diagnosis \\
Provides phase-specific communication strategies and sample language
};

\node[box, fill=fillCore, below=of navigation] (rag) {
\textbf{Retrieval-Augmented Generation (Objective 2)} \\
Grounds LLM responses in Ariadne Labs Essential Communications Toolkit \\
Ensures source-faithful, evidence-based guidance
};

\node[box, fill=fillCore, below=of rag] (stuck) {
\textbf{Stuck Points Framework (Objective 4)} \\
Identifies and addresses relational communication obstacles \\
Strategies: Acknowledge, Get Curious, Summarize and Plan
};

% Partial contribution (lighter green fill)
\node[box, fill=fillPartial, below=of stuck] (llm) {
\textbf{LLM-powered conversational interface} \\
Real-time contextual guidance during consultations \\
Dual-mode architecture for multi-provider comparison
};

% Context / not contributed (grey fill)
\node[box, fill=fillContext, below=of llm] (eval) {
\textbf{Evaluation and validation} \\
Evidence grounding \quad • \quad Source fidelity \quad • \quad Clinical relevance \\
{\color{accentGreen}\textbf{Framework integration and user-centered design}}
};

% Arrows
\draw[arrow] (patient) -- (navigation);
\draw[arrow] (navigation) -- (rag);
\draw[arrow] (rag) -- (stuck);
\draw[arrow] (stuck) -- (llm);
\draw[arrow] (llm) -- (eval);

% Legend (fill-based semantics)
\node[box, fill=white, below=10mm of eval] (legend) {
\textbf{Legend} \\[2pt]
\begin{tabular}{@{}p{0.30\textwidth} p{0.66\textwidth}@{}}
\textbf{\colorbox{fillCore}{\strut\hspace{14pt}}} &
Core contribution: Navigation Map, RAG pipeline, Stuck Points Framework. \\
\textbf{\colorbox{fillPartial}{\strut\hspace{14pt}}} &
Partial contribution: LLM interface and dual-mode architecture. \\
\textbf{\colorbox{fillContext}{\strut}} &
Context or evaluation components. \\
\textbf{\colorbox{fillExisting}{\strut\hspace{14pt}}} &
Pre-existing clinical setting and patient context. \\
\end{tabular}
};

\end{tikzpicture}

\caption{
End-to-end overview of an LLM-guided dementia consultation system.
Box fill indicates contribution level; green text marks key thesis contributions.
}
\label{fig:graphical_abstract_overview}
\end{figure}
